% =====================================================================
%  TRABAJOS FUTUROS (versión por objetivos)
%
%  Corrección del director del 11 de noviembre de 2026: las líneas que había
%  eran mejoras técnicas de lo hecho, así que van en un párrafo y una tabla; los
%  trabajos futuros propiamente dichos salen del estado del arte.
% =====================================================================
\ustacapitulocontinuo{Trabajos futuros}
\addcontentsline{toc}{chapter}{TRABAJOS FUTUROS}

Las limitaciones de la plataforma y de la operación que mostraron las pruebas se
pueden atender con las mejoras de la tabla~\ref{tab:mejoras}, que no cambian la
arquitectura del sistema.

\begin{table}[H]
  \centering\footnotesize\setstretch{1.0}
  \caption{Mejoras técnicas del sistema y la limitación que atiende cada una.}
  \label{tab:mejoras}
  \begin{tabular}{@{}>{\raggedright\arraybackslash}p{0.17\textwidth}>{\raggedright\arraybackslash}p{0.33\textwidth}>{\raggedright\arraybackslash}p{0.42\textwidth}@{}}
    \toprule
    \textbf{Mejora} & \textbf{Limitación que atiende} & \textbf{Cómo} \\ \midrule
    Navegación más precisa & El agente no se mueve por debajo de 0,40~m/s y solo mide su avance donde el pasillo tiene estructura & Calibrar el variador o cambiar la relación de transmisión, y medir el avance con codificadores en las ruedas u odometría visual con las cámaras que el agente ya lleva; la tolerancia de llegada podría volver a 0,25~m \\ \addlinespace
    Más capacidad de cómputo & La tarjeta, de dos núcleos, queda entre el 54~\% y el 59~\% de uso en reposo con la navegación completa & Adaptar la plataforma a una Raspberry Pi~4 o~5, de cuatro núcleos; la comunidad de AWS DeepRacer ya publica el software para la Raspberry Pi~4 \\ \addlinespace
    Evasión de obstáculos & Las pruebas se hicieron con los pasillos despejados & Esquivar personas y objetos en movimiento, y no solo detenerse ante ellos, y probarlo con personas que caminan por el pasillo antes de operar en horario de clases \\ \addlinespace
    Controlador predictivo & Regulated Pure Pursuit anula la curvatura cuando una inversión de marcha queda cerca de la meta & Usar MPPI, que incluye Nav2 y admite cinemática tipo automóvil, para devolver la media vuelta a Nav2 y afinar la aproximación final; exige más cómputo \\ \addlinespace
    Interacción por voz & El usuario elige el origen y el destino en listas & Reconocimiento de voz a texto y un modelo de lenguaje en el computador del coordinador, que solo traduzca la petición a un origen y un destino del catálogo \\ \bottomrule
  \end{tabular}
\end{table}

\paragraph{Más de dos pisos.} Con más pisos, una misión sigue teniendo un solo
relevo, entre el agente del piso de origen y el del piso de destino. Lo que cambia es
que varias solicitudes a la vez compiten por los mismos agentes, y decidir cuál atiende
cada tramo y en qué orden puede requerir procesos de optimización multiagente. Una
alternativa que hoy se encuentra en el estado del arte es el
aprendizaje por refuerzo multiagente (MARL), que \textcite{xu2026modular} integran en
un marco modular sobre ROS para asignar tareas entre varios robots y cuyas
aplicaciones en sistemas multirrobot revisan \textcite{orr2023multiagent}.

\paragraph{Varios agentes en un mismo piso.} Así mismo, una oportunidad identificada
es la inclusión de MARL para coordinar varios agentes que operen sobre un mismo piso.
Con dos o más agentes elegibles por tarea, la asignación deja de ser una regla fija y
hacen falta una cola de solicitudes y la noción de agente ocupado; además de los
mecanismos de mercado, como la subasta \autocite{quinton2023market}, MARL permitiría
aprender esa coordinación a partir de la operación. Es la extensión natural para
edificios con más afluencia, y obliga a medir de nuevo el costo del descubrimiento en
la red, que crece con el número de procesos.

\paragraph{Usuarios que no pueden usar la escalera.} El relevo supone que el usuario
cambia de piso por la escalera. Para quien no puede usarla, el punto de transferencia
tendría que estar en el ascensor, y la elección entre escalera y ascensor, que
\textcite{luo2025elevator} resuelven para el robot con aprendizaje por refuerzo
jerárquico, pasaría a depender del usuario.
